\chapter{An operational protocol}\label{ch:protocol}

Experienced release committees often regard a written protocol with suspicion, and
the suspicion has a fair basis. The safety of a release comes from the protection
applied to it and from the judgement of the people who know the data and its users.
Committees made sound decisions long before anyone asked them to fill in a record
with nine elements. A protocol adds forms, sign-offs and waiting time, and none of
that changes a single cell of the file. On this view a protocol is paperwork that
slows releases without making them safer.

The view fails on the evidence of the method's own audit. Suppose a disclosure
analyst runs an honest audit of an attack whose true risk sits exactly at the
tolerance. The audit certifies it falsely in 3.8 per cent of runs, inside its stated
error level of 5 per cent. Now let the analyst run the same audit ten times with
different seeds and report the first that passes. The false certification rate
becomes 32 per cent, and with a hundred tries 98 per cent.
% source: paper/publication_versions/table_audit_process.tex (0.038, 0.322, 0.979)
Nothing in the file changed, and no single audit was computed wrongly. What changed
is that the choice of seed was made after the outcomes could be seen. The control
that restores the rate of 3.8 per cent is a matter of order: the targets are fixed
before any outcome exists, and every audit is recorded where it cannot be dropped.
That is a protocol. The error statement that makes a risk bound worth computing
holds only when certain things are written down before certain other things happen.

This chapter asks what an institution must write down, in what order and by whom,
for a release decision to carry the error statement the method promises.

\section{The order is the control}

Every guarantee in the earlier chapters has the same shape. A bound computed from
simulated attacks lies below the true risk with probability at most $\alpha$, and so
a release whose true risk exceeds the tolerance passes with probability at most
$\alpha$. The probability is taken over the attacks, with everything else held fixed:
the tolerance, the declared adversary, the target set, the setting under test and
the number of attacks. If any of these is chosen after the outcomes are seen, the
probability statement is about a different experiment, and in the worst case it says
nothing.

Proposition~\ref{c16:prop-order} states this for the two choices an institution is
most tempted to make late: the tolerance itself, and which of several audits to
report.

\begin{proposition}[The order of declaration]\label{c16:prop-order}
Let $\overline{R}$ be an upper bound with $\Prob(\overline{R} < R) \le \alpha$ for the
true risk $R$ of a declared attack.
\begin{enumerate}
\item[(i)] If the tolerance $\tau$ is fixed before the attacks are run and
$R > \tau$, the release passes with probability at most $\alpha$.
\item[(ii)] If the tolerance may be set after the attacks, as any value in a range
that contains $\overline{R}$, then a rule that sets $\tau = \overline{R}$ passes
every release, whatever its risk.
\item[(iii)] If the operator can run $N$ independent audits of the same attack and
report one that passes, and a single audit passes with probability $f$, the
reported audit passes with probability $1 - (1-f)^N$, which tends to one as $N$
grows whenever $f > 0$.
\end{enumerate}
\end{proposition}

Part (iii) produces the figures of the opening paragraph: with $f = 0.038$, ten tries
give $1 - 0.962^{10} = 0.32$ and a hundred give 0.98. Table~\ref{c16:tab-process}
reports the full measurement from the repository, for an attack at the tolerance of
0.12 and for one slightly above it at 0.14, and adds a third way of breaking the order:
an operator who chooses the bet of a sequential audit after seeing the outcome
certifies an unsafe attack in every run. With the controls in place, each row returns
to the rate of a single honest audit.

\begin{table}[ht]
\centering
\small
\caption{Attacks on the audit process: the share of runs in which an attack at or
above the tolerance of 0.12 is certified, with one coordinate at the full level
$\alpha = 0.05$. ``Tries'' is the number of seeds or unregistered shadow audits the
operator can choose among. Without the control, the rate for several tries is
$1 - (1-f)^{\text{tries}}$ for the single-try rate $f$; with it, the operator has one
stream of targets fixed in advance.}
\label{c16:tab-process}
% source: paper/publication_versions/table_audit_process.tex; scripts/audit_process_attacks.py
\begin{tabular}{lrrrr}
\toprule
Attack on the audit & True risk & Tries & Without control & With control \\
\midrule
Seed grinding or shadow audits & 0.12 & 1 & 0.038 & 0.038 \\
Seed grinding or shadow audits & 0.12 & 10 & 0.322 & 0.038 \\
Seed grinding or shadow audits & 0.12 & 100 & 0.979 & 0.038 \\
Seed grinding or shadow audits & 0.14 & 10 & 0.066 & 0.007 \\
Seed grinding or shadow audits & 0.14 & 100 & 0.495 & 0.007 \\
Bet chosen after the outcome & 0.14 & 1 & 1.000 & 0.007 \\
Bet chosen after the outcome & 0.20 & 1 & 1.000 & 0.000 \\
\bottomrule
\end{tabular}
\end{table}

The table does not describe dishonest analysts. Each row describes something an
analyst under a deadline might do without intending harm: rerun an audit that
failed because the seed looked unlucky, try a few settings in a scratch copy before
registering the one that works, or adjust the audit's design after a bad early
result. A protocol that fixes the order of work removes these choices, and with them
the inflation they cause. The paperwork is the means by which the order can be
checked afterwards.

Grace, from the Bank's internal audit function, reads decision records after the
fact, and she reads them for order. Her first question on any record is not whether
the risk figure is small. It is whether the tolerance was signed before the first
attack outcome existed, and whether every audit run on the register appears in the
registry. A record that cannot answer both questions, she reports, carries no error
statement at all, however careful its arithmetic.

\section{Four roles and three lines}

The protocol needs four roles, and the first and third should not be held by the
same person. The \emph{release owner} is a named senior officer, at the Bank the
chair of the statistics release committee, who sets the tolerance, signs the
decision and owns the residual risk. The \emph{data owner} is the business area that
collects the source register and knows which products are drawn from it, how often
and under what legal basis. The \emph{disclosure analyst}, Naledi at the Bank,
declares the scenarios, runs the attacks and computes the bounds; she owns the
correctness of the calculation and not the decision. The \emph{reviewer}, outside
the production line in a risk, audit or methodology function, checks that the
scenarios are plausible and that the falsification tests of
Chapter~\ref{ch:falsification} were run.

The roles map onto the three lines of defence that central banks already use for
operational risk. The analyst and the data owner sit in the first line, the
reviewer in the second, and internal audit, which checks the filed record and the
controls on the audit process, in the third. The separation is the point. The
person who computes the bound should not be the person who decides whether the
residual risk is acceptable, and neither should be the person who checks afterwards
that the work was done in the stated order. Chapter~\ref{ch:thresholds} sets out the
legal basis for this division and the elements of the threshold the owner signs.

\section{Seven steps}

Figure~\ref{c16:fig-steps} shows the seven steps as a decision process, with the
role that leads each step. Two loops leave the main line. The dotted loop is the
re-assessment at each new vintage or at expiry, which begins again at the planning
step and draws on what remains of the programme's error budget. The dashed loop is a
change of product: when no setting meets the tolerance, the institution moves to a
different product, such as aggregate tables or access in a secure setting, and
assesses it from the start.

\begin{figure}[ht]
\centering
\begin{tikzpicture}[font=\small,
  stp/.style={draw,align=center,text width=6.4cm,minimum height=0.62cm},
  an/.style={stp},
  ow/.style={stp,fill=gray!20},
  rv/.style={stp,pattern=north east lines,pattern color=gray!55}]
  \node[an] (s1) at (0,0) {1. List the programme, not the release};
  \node[an] (s2) at (0,-0.95) {2. Declare the scenarios};
  \node[ow] (s3) at (0,-1.9) {3. Set the tolerance, owner, $\alpha$ and expiry};
  \node[an] (s4) at (0,-2.85) {4. Plan the measurement};
  \node[an] (s5) at (0,-3.8) {5. Measure per release and pooled};
  \node[rv] (s6) at (0,-4.75) {6. Test the model behind the number};
  \node[ow] (s7) at (0,-5.7) {7. Decide, record and re-assess};
  \foreach \a/\b in {s1/s2,s2/s3,s3/s4,s4/s5,s5/s6,s6/s7} \draw[->] (\a) -- (\b);
  \draw[->,densely dotted,thick] (s7.west) -- ++(-0.7,0) |- node[pos=0.25,left,align=right,font=\scriptsize]{new vintage\\or expiry} (s4.west);
  \draw[->,dashed,thick] (s7.east) -- ++(0.7,0) |- node[pos=0.25,right,align=left,font=\scriptsize]{no setting passes:\\change the product} (s1.east);
  \node[draw,minimum width=0.5cm,minimum height=0.3cm] at (-3.4,-6.75) {};
  \node[anchor=west,font=\footnotesize] at (-3.1,-6.75) {analyst};
  \node[draw,fill=gray!20,minimum width=0.5cm,minimum height=0.3cm] at (-1.5,-6.75) {};
  \node[anchor=west,font=\footnotesize] at (-1.2,-6.75) {release owner};
  \node[draw,pattern=north east lines,pattern color=gray!55,minimum width=0.5cm,minimum height=0.3cm] at (1.4,-6.75) {};
  \node[anchor=west,font=\footnotesize] at (1.7,-6.75) {reviewer};
\end{tikzpicture}
\caption{The seven steps of the protocol and the role that leads each one. Plain
boxes are led by the disclosure analyst, grey boxes by the release owner and hatched
boxes by the reviewer. The dotted arrow is the re-assessment at a new vintage or at
expiry; the dashed arrow is a change of product, which needs a new assessment.}
\label{c16:fig-steps}
\end{figure}

\paragraph{Step 1. List the programme, not the release.} The analyst starts from the
source register and not from the product about to be published. Every product drawn
from the register over the planning horizon is listed: monthly or quarterly
aggregates, supervisory publications, research extracts, data shared with other
authorities and answers to ad hoc requests. For each, the inventory records its
frequency, the population it covers, the variables it publishes and whether any
identifier or stable ordering persists from one vintage to the next. A persistent
pseudonym, a stable row order or a rare combination of fields that seldom change lets
an adversary match records across vintages, which is the condition of
Proposition~\ref{prop:nopool}. The step ends with a declared horizon $m$, the number
of vintages the assessment covers.

\paragraph{Step 2. Declare the scenarios.} For each product, and for each
combination of products an outsider could hold together, the analyst writes down who
the adversary is, what it knows, which records it targets and what harm follows if it
succeeds. For a pooled scenario the declaration says whether the adversary knows in
which vintages the target is present, and why, because on the credit register that
single piece of knowledge decides whether pooling adds anything
(Proposition~\ref{prop:nohistory}). The declaration should be specific enough that
someone can disagree with it. Chapter~\ref{ch:worksheets} gives starting declarations
for six release types.

\paragraph{Step 3. Set the tolerance.} For each harm class the release owner sets a
tolerance, an error level, an expiry date and the horizon of Step~1, and records the
legal basis, the parties the tolerance protects and the access route. The tolerance
can be stated directly, as the largest probability of successful exposure the
institution accepts for that harm, or indirectly, as a relative weight on harm
against lost utility from which the tolerance follows
(Chapter~\ref{ch:tolerance}). Either way it is written down before any risk is
measured. Where no tolerance has been agreed for a harm, the record says so and no
decision is made on that harm.

\paragraph{Step 4. Plan the measurement.} Before running anything, the analyst fixes
how many attacks each scenario needs. With no observed success, certifying a
tolerance $\tau$ at $\alpha = 0.05$ needs about $3/\tau$ independent attacks
\citep{hanley1983}; exactly, 299 for a tolerance of 0.01 and 2{,}995 for 0.001.
% source: paper/generated/table_certification.tex
If a floor is to be chosen from a grid, the order in which the floors are tested is
fixed now (Chapter~\ref{ch:selection}). If the release must pass every one of a fixed
set of scenarios, each may be assessed at the full level, since that is an
intersection--union test. If scenarios may be added later, for example because a
reviewer proposes a stronger adversary, the plan reserves part of the error budget
for them (Chapter~\ref{ch:monitoring}).

\paragraph{Step 5. Measure, per release and pooled.} The analyst runs the declared
attacks on the release as it would be published, on fresh targets, and computes the
exact upper bound for each scenario. The pooled attack is then run over the
programme of Step~1 with the same auxiliary information. Pooled attacks that add
knowledge, such as more readings or the target's history across vintages, are
reported separately, because a committee needs to know which adversary it is
protecting against. Where a level is chosen from a grid, it is audited again on fresh
targets. The step reports what the chosen level keeps, in records and in volume.

\paragraph{Step 6. Test the model behind the number.} Before the bound is used, the
reviewer confirms that the six falsification tests of
Chapter~\ref{ch:falsification} were run against bands fixed in advance. A failed
release-level test makes the assessment invalid, not failed: the institution cannot
release on it and must repair the model first.

\paragraph{Step 7. Decide, record and re-assess.} The release passes a scenario only
if the upper bound, not the point estimate, is at or below the tolerance, and it
passes overall only if it passes every declared scenario, pooled ones included. When
the programme fails, the release fails, even if it passes on its own. The owner signs
the decision record. Re-assessment is due when the programme is about to pass its
horizon, when a new product is drawn from the register, when a new public source
could serve as auxiliary information, when the mechanism changes and at expiry.

The long form of the paper lists six stages where this list has seven. The
difference is one of presentation: the long form folds the inventory of Step~1 into
its first declaration stage, and the planning of Step~4 into its computation stage.
% source: paper/bis_ifc_full.tex sec:protocol (six stages); paper/apply_section.tex sec:apply-steps (seven steps)
The seven-step version is the one an institution can staff, because each step has
one output and one lead.

When Grace audited the record for Tomás's third request, she checked the steps
against their dates. The inventory listed the three extracts and the quarterly
aggregates drawn from the same register. The scenario sheet was dated a week before
the tolerance was signed, and the tolerance was signed two days before the registry
showed its first attack. She noted one finding: the reviewer's sign-off on the
falsification tests was dated after the decision. The tests had been run on time;
the sign-off was late. She recorded it as a control weakness, because a late
sign-off cannot show that the decision waited for it.

\section{A filled decision record for the credit register}

Table~\ref{c16:tab-record} fills the record for the synthetic credit register of
Chapter~\ref{ch:credit}, for the first vintage and for the programme of six. It
rebuilds the decision record of the long form of the paper and uses only figures that
the long form reports from the repository's runs.
% source: paper/bis_ifc_full.tex tab:record
The register is synthetic, so the record shows how the method fills a record, not
what any central bank would find. The tolerance of 0.12 is illustrative, as are the
owner, the dates and the legal basis.

The declarations common to every row sit outside the table, because they are the
same for each harm. The legal basis is Council Regulation (EC) No 2533/98,
Article~8, as amended by Council Regulation (EC) No 951/2009 \citep{ec2533,ec951},
given here for illustration. The protected parties are the borrowers, the reporting
agents and the natural persons behind borrowers, such as sole traders. The access
route is a scientific-use file for accredited researchers. The error level is
$\alpha = 0.05$. The owner is the chair of the statistics release committee. The
record expires one year after the assessment, and it is re-assessed earlier on a
vintage beyond the horizon $m = 6$, on a new auxiliary source available to the
adversary, on a change of mechanism or on a change of access route. The floor of 150
was chosen by fixed-sequence testing from 300 down and confirmed on a second,
independent draw of targets.

\begin{table}[p]
\centering
\caption{Decision record for the synthetic credit register, on the first vintage as
released and on the programme of six vintages, tolerance 0.12 where one is set.
Declarations common to every row are given in the text. The bound is the exact
one-sided 95 per cent upper bound from $B$ simulated attacks; where a row reports a
census of released cells instead of a simulation, the measure is named.}
\label{c16:tab-record}
% source: paper/bis_ifc_full.tex tab:record (every figure in this table)
\footnotesize
\setlength{\tabcolsep}{3pt}
\begin{tabular}{@{}>{\raggedright\arraybackslash}p{1.9cm}>{\raggedright\arraybackslash}p{2.0cm}>{\raggedright\arraybackslash}p{2.4cm}>{\raggedright\arraybackslash}p{0.7cm}>{\raggedright\arraybackslash}p{2.6cm}>{\raggedright\arraybackslash}p{1.8cm}>{\raggedright\arraybackslash}p{2.2cm}@{}}
\toprule
Harm & Target set & Adversary & $\tau$ & Evidence & Setting & Decision \\
\midrule
Identification, one release & Released loans, first vintage & Linkage attack with
one exposure reading & 0.12 & At $k_{\min} = 150$: confirmation 0.0135, bound
0.017, $B = 4{,}000$; this row alone would need no floor & $k_{\min} = 150$ & Pass \\
\addlinespace
Identification, programme & Loans over six vintages, $m = 6$ & Pooled attack, one
reading per live vintage, knows the loan's history & 0.12 & Floor 150 by fixed
sequence; confirmation 0.101, bound 0.109, $B = 4{,}000$ fresh targets; keeps 20\%
of loans, 4\% of exposure & $k_{\min} = 150$ & Pass on risk; owner decides whether
the file is still worth releasing; re-assess before vintage 7 \\
\addlinespace
Identification, every size class & Each size class, largest firms on a stratified
draw & As in the second row & 0.12 & At $k_{\min} = 150$ the worst class bound is
0.184; no floor up to 200 passes & None found & Fail; offer aggregates or secure
access \\
\addlinespace
Identification by an insider & Released loans & A party to the loan, log exposure
error 0.01, knows the history & Not set & Success 0.772 at $k_{\min} = 1$, 0.416 at
$k_{\min} = 300$ & Not governed by $k_{\min}$ & Fails at every floor once a tolerance
is set; restrict access \\
\addlinespace
Relationship allegation & Released borrowers & Names the largest lender in the key
class & Not set & Predictive value 0.328 at $k_{\min} = 1$, 0.313 at $k_{\min} = 20$
& Not governed by $k_{\min}$ & No decision until $\tau$ is set \\
\addlinespace
Arrears status & Large borrowers & Reads status off a homogeneous key class & Not
set & Homogeneity share 0.119 at $k_{\min} = 20$ & Attribute rule needed & No
decision until $\tau$ is set \\
\addlinespace
Contribution from totals & 814 published cells & A participating lender, by
subtraction & -- & 5.9\% of cells recoverable & Three-contributor rule; dominance
rule not assessed & Rule-based; not a probability \\
\addlinespace
Differencing & Consecutive cell pairs & Claims a single move on a unit change & Not
set & A unit change is a single move in 65\% of cases at $c_{\min} = 1$, 56\% at
$c_{\min} = 10$ & $c_{\min} = 10$ & No decision until $\tau$ is set \\
\bottomrule
\end{tabular}
\end{table}

The record is worth reading row by row, because its value lies in what the rows
disagree about. The first row is the assessment a committee would have made under
per-release review: one release, the outside adversary with one reading of exposure,
and a bound of 0.017 at the chosen floor. On its own this row would need no
suppression at all. The second row is the programme. The record declares the full
outside adversary, who pools six vintages, reads exposure in each, and knows in which
of them the target's loan was live, and against that adversary the floor of 150
passes with a confirmed bound of 0.109. The row then states what the floor costs, a
fifth of the loans and 4 per cent of the exposure, because the owner has to weigh
that cost and cannot do so from a risk figure alone.

The third row fails. When every size class must meet the tolerance, with the largest
firms drawn on purpose, the worst class bound at the floor of 150 is 0.184 and no
floor up to 200 passes. The fourth row names the insider, Pieter's position, whose
success is 0.772 with no suppression and 0.416 at a floor of 300. No floor governs
him; access does. The next two rows, the relationship allegation and arrears status,
have no tolerance, because the institution has not decided how much of those harms it
will accept. The record says so and makes no decision on them. The seventh row is
governed by a rule and not by a probability: a participating lender can recover 5.9
per cent of the 814 published cells by subtraction, and the three-contributor rule
answers it. The last row shows a differencing finding with a floor applied and no
tolerance yet.

Read together, the rows decide the file. One row passes, one passes at a cost the
owner must accept, one fails at every floor tested, one would fail at every floor
once a tolerance is set, three cannot reach a decision, and one is governed by a rule.
The decision on the file as a whole follows from those rows and not from the one that
passes. The long form draws the conclusion for this synthetic register: against an
adversary who knows loan histories, a loan-level file is the wrong product, and
aggregates or secure access are the options a committee can defend.

The record also shows what per-release review would have hidden. The per-release
reading of the rule releases all 27{,}167 loans of the first vintage; the programme
reading against the full adversary releases 5{,}553, with 4 per cent of the exposure.
% source: paper/bis_ifc_full.tex sec:disc (27,167 and 5,553); paper/generated/numbers.tex runevalperreleasekept
Whether the adversary knows the loan's history moves the programme floor from 1 to
75, and whether it also holds a reading in every vintage moves it from 75 to 150.
% source: paper/generated/table_robustness.tex
These are the differences between declarations, and a record is the only place where
a committee can see them side by side.

The chair's first reaction to the record was that it made the decision harder. A
single pass would have been easier to sign. Naledi's answer was that the single pass
existed, in the first row, and that it described an adversary the committee did not
believe in. The chair signed the record with the third row's remedy: Tomás would
work with the loan-level file inside the Bank's secure environment, and the extract
that left the building would be aggregate.

\section{When a release fails}

A failed scenario does not end the matter. It tells the committee what to change,
and the record states which remedy was chosen and why. The remedies fall into five
groups, in roughly increasing order of what they cost the user.

The first is to change the key: coarsen the published breakdown, or withdraw the
field the adversary uses most. On the synthetic register the exposure reading, not
the key, carries the attack; withdrawing it lowers identification from 0.081 to
0.025. % source: paper/generated/numbers.tex runevalablationfull 0.081, runevalablationnoexposure 0.025
The second is to raise the floor, choosing the level by fixed-sequence testing on the
upper bound of the pooled scenario and confirming it on fresh targets. The third is to
break the link between vintages by rotating pseudonyms, changing the row order or
publishing some products less often. Rotation controls linkability and does not make
a record anonymous in the legal sense, and the pseudonym should be a keyed hash under
a key held apart from the release team. It helps only if the published fields cannot
relink the records: on the synthetic register the key and the exact exposure relink
every loan from one vintage to the next, and with exposure rounded to one significant
figure about one loan in ten (Chapter~\ref{ch:programmes}). The fourth is to add noise
under a formal mechanism and account for its budget over the horizon, remembering
that the guarantee bounds posterior odds and does not name a linkage probability
(Chapter~\ref{ch:dp}). The fifth is to change the access route or the product: move
the file from public or scientific-use release to a secure research environment,
where access controls and output checking govern the risk \citep{desai2016}, or
publish aggregates in place of loan-level records.

Shortening the horizon is also a remedy, but only if the institution is prepared to
stop the series or re-assess it at the new horizon. A horizon shortened on paper
while the series continues is no remedy at all.

\section{Starting small}

An institution does not need to apply the protocol to every product at once. A
workable start is one register and one product family over one quarter. Choose the
register whose products are published most often. Complete the programme inventory
for it and agree a horizon. Agree one tolerance, for identification, with a named
owner. Run one per-release attack and one pooled attack, with the attack counts of
Step~4. File one decision record and present it to the release committee beside the
current method, so that the committee can see where the two disagree.

The pilot's purpose is that comparison. A committee that sees its current method and
the protocol reach the same answer has learnt that its method was sound for that
product and can say why. A committee that sees them disagree has found the case that
the protocol was built for, and has found it on one product, before it affects
many. The pipeline that produces this book runs the attacks and writes the record; it
needs, for each product, the published file or its specification, the key fields, the
auxiliary variables the adversary is declared to hold and the tolerance. Its
demographic path has run on two public files (Chapter~\ref{ch:public}); its credit
path has so far run only on the synthetic register.

\section{What the method adds to the security of a release system}

A statistical release system can be assessed in the terms used for any other
information system, and readers in information security, operational risk and
internal audit will find the method easier to place in those terms. This section
restates it so, following the long form. It does not claim that the method meets, or
is certified under, any security standard.

The asset is the confidential register: its units and their attributes. Disclosure
review faces two adversaries. The first, $\mathcal{A}_{\mathrm{data}}$, holds the
releases and the declared auxiliary information and tries to identify a unit; on the
credit register the record declares three outside classes of it and reports the
insider beside them. The second, $\mathcal{A}_{\mathrm{op}}$, is whoever runs the
audit and wants a release to pass, through carelessness, deadline pressure or
intent. Per-release review has a model of neither. The method gives a claim about the
first that holds only if the second is contained, so both belong in the record.

The claim is this. For each declared adversary and harm class, with probability at
least $1 - \alpha$, no release is certified whose declared attack succeeds at a rate
at or above the tolerance. The probability is over the simulated attacks, given the
register and the declared attack. It holds against $\mathcal{A}_{\mathrm{op}}$ only if
five things are trusted: the registry that stores the audits and the error budget,
one per programme; the randomness from which targets are drawn, which the operator
cannot predict when it registers an audit; the attack code, fixed by a digest when it
is registered; the list of attacks declared for each release, fixed before its first
outcome; and a time source the operator does not control. An operator who chooses the
seed knows every outcome in advance, and then no error statement holds, which is the
content of Table~\ref{c16:tab-process}.

The repository implements matching controls in its adaptive audit module. Targets are
derived from a secret committed at registration and a later round of a public
randomness beacon, so that neither the operator nor the beacon alone chooses them.
Each bet is recorded in a hash-chained ledger before the attack runs, and its outcome
after. The ledger's entries cannot be changed, dropped or reordered without breaking
the chain or changing its head. With the controls in place, and the registry trusted,
the rates in the table return to those of an honest single audit.

Against the outside adversary the method adds four properties to per-release review.
It assesses the pooled attack over the declared horizon, which per-release figures
cannot bound (Proposition~\ref{prop:nopool}). It compares an upper bound with the
tolerance, so the chance of false assurance is at most $\alpha$, where a point
estimate passes a release at the boundary about half the time; and when the evidence
is watched as it arrives, it uses a sequence that stays valid, where twenty looks at a
fixed-sample bound put it below the true risk in 21 per cent of runs.
% source: paper/generated/numbers.tex runseqfixederror 21, runseqanytimeerror 0, runseqlooks 20
It fails closed: an attack that has not been audited is given a bound of one, and the
pipeline returns a fail when the pooled bound exceeds the tolerance even if each
release passes on its own. And a new attack proposed by a red team enters with its
own share of the error budget, without voiding certificates already issued, if budget
was reserved for it. Against the operator, the method makes audit shopping, the
repetition of an audit or the scanning of settings until one passes, either impossible
or paid for out of the error budget. That holds only under complete mediation
\citep{saltzer1975}: every attack run on the confidential register passes through the
registry, and development work uses a separate extract.

What the method does not give should be stated as plainly. The implementation is a
library, not a security boundary: read-only fields in Python are not a security
control, since code in the same process can change them. The SHA-256 hashes a run
writes show that a file has changed since the run, but they carry no key or trusted
time stamp, and the ledger's chain detects changes only if its head is published or
time-stamped outside the operator's control. A deployment that must hold against the
operator needs the registry as a separate service with its own storage, signing keys
and time stamps, which the repository specifies and has not built.
% source: docs/DEPLOYMENT_AND_PENTEST_2026_10_23.md ("The registry service is specified in the deployment paper and not built")
It needs separation of duties, in which an independent team writes and signs the
attack code, a data custodian holds the register and the randomness commitment, and
the owner signs the tolerance before the first outcome. And it needs verification by
internal audit inside the secure environment, since the log of which targets were
identified is itself confidential. The method gives no guarantee against attacks that
were not registered, and no differential privacy guarantee.

In the language of security practice, the declared scenario is a threat model
\citep{shostack2014}: asset, adversary, capability and harm written down before
testing. The attack simulation is red-team testing with a statistical stopping rule.
The decision record is an assurance case: a claim, the evidence for it and the
argument that links them, with a named owner and an expiry date. The mapping helps an
institution place the method within its risk and audit processes. It is not a claim
of conformity.

For Grace, this section is the part of the method she can test. She cannot judge
whether the declared adversary is realistic; that is the reviewer's work. She can
check that the registry holds every audit, that the targets were derived from a
commitment made before the beacon round, that the ledger's head matches the copy
published at registration, and that the tolerance's signature predates the first
entry. Each of those checks either passes or fails on the documents.

\section{What the protocol does not decide}

The protocol makes a decision reproducible. It does not make the decision for the
institution. It cannot say whether a declared adversary is realistic, whether a
tolerance is legally or ethically right, or whether a release is worth its risk.
Those remain judgements for the release owner. The repository's own register of
claims puts the protocol's benefit in the same terms: that it improves auditability is
an institutional argument, supported by a machine-readable record, generated tables
and hashes that show it can be done, and it would be refuted by comparative evidence
that it does not help or costs too much. % source: docs/CLAIMS_EVIDENCE.md last row
No such comparison has been run.

The answer to the chapter's question can now be given plainly. For a release
decision to carry the error statement the method promises, the institution must
write down, before any attack outcome exists, the programme and its horizon, the
declared adversaries, the tolerance with its owner, error level and expiry, and the
plan of measurement, including the order in which settings are tested. The analyst
then measures and the reviewer tests the model, each recorded in a registry that
holds every audit. The owner decides from the full record, and internal audit checks
the order afterwards. The writing is not overhead on a safe release. Without it the
probability statement attached to the bound does not hold, and with ten unrecorded
tries an audit at the tolerance passes nearly a third of the time.

\begin{inpractice}[In practice: running the protocol]
\setlength{\emergencystretch}{3em}%
Assign the four roles to named people, and keep the analyst and the owner apart.
Date every output: the inventory, each scenario sheet, the signed tolerance, the
measurement plan, every registered audit, the falsification results and the signed
decision. The tolerance's signature must predate the first audit outcome, and every
attack run on the confidential register must appear in the registry, including
failed and abandoned runs. File a record with one row per harm, including the harms
for which no tolerance has been set. When a release fails, record the remedy and
assess the changed product from Step~1. Start with one register and one product
family, and compare the protocol's record with the current method before extending
it.

\smallskip\noindent\textit{In the software.} \texttt{privacy\_\allowbreak{}risk.\allowbreak{}pipeline.\allowbreak{}gate}
writes the decision record (\texttt{decision.\allowbreak{}json}) with the release and programme
outcomes, and \texttt{governance.\allowbreak{}decide} applies the tolerance and the expiry date.
The audit controls are in \texttt{privacy\_\allowbreak{}risk.\allowbreak{}adaptive\_\allowbreak{}audit}, and
\texttt{scripts/\allowbreak{}audit\_\allowbreak{}process\_\allowbreak{}attacks.\allowbreak{}py} produces Table~\ref{c16:tab-process}. The
governance record the code writes holds six of the nine elements of the threshold. The other three, the legal basis with
the protected parties, the access route and the horizon, have to be added by hand
(Appendix~\ref{app:software}, Chapter~\ref{ch:thresholds}).
\begin{lstlisting}
@dataclass(frozen=True)
class GovernanceThreshold:
    tolerance: float
    owner: str
    expiry_date: str
    adversary: str
    harm_class: str
    alpha: float
\end{lstlisting}
\end{inpractice}

\section{Proofs}

\begin{proof}[Proof of Proposition~\ref{c16:prop-order}]
(i) With $\tau$ fixed before the attacks, the pass event $\{\overline{R} \le \tau\}$ is
contained in $\{\overline{R} < R\}$ whenever $R > \tau$, and the latter has probability
at most $\alpha$.

(ii) If $\tau$ may take any value in a range containing $\overline{R}$, the rule
$\tau = \overline{R}$ is admissible, and under it $\overline{R} \le \tau$ holds on every
outcome, so the release passes with probability one whatever $R$ is.

(iii) Let $P_1,\dots,P_N$ be the pass indicators of $N$ independent audits, each with
$\Prob(P_i = 1) = f$. The operator can report a passing audit exactly when at least
one passes, which has probability $1 - \prod_i \Prob(P_i = 0) = 1 - (1-f)^N$. For
$0 < f \le 1$, $(1-f)^N \to 0$ as $N \to \infty$.
\end{proof}
